% ================================================
% 文件: 01_电台设备.tex
% 章节: 第14章 电台设备
% 所属部分: 系统与设备
% 版本: v1.0
% ================================================
% 本章目录结构（树状）:
% \chapter{电台设备}
% ├── \section{电台设备概述}
% ├── \section{电台组成与原理}
% ├── \section{电台类型与分类}
% ├── \section{电台技术指标}
% └── \section{电台应用与维护}
% ================================================

\chapter{电台设备}
\label{chap:radio_station}

\section{电台设备概述}
\label{sec:radio_station_intro}

电台设备是用于无线电通信的完整系统，包括发射机、接收机、天线、电
源等部分。电台广泛应用于军事、民用通信、业余无线电等领域，既可以
是固定安装的基地台，也可以是车载、船载或便携的机动平台。与单一发
射机或接收机相比，电台强调"收发一体"与整机协同工作，是通信网中
的最小节点单元。

电台区别于独立收发信机的主要特点体现在四个方面。其一是\text
bf{双向通信}：同时具备发射和接收功能，通过收/发切换或双工器
共用一副天线；其二是\textbf{可移动性}：便携式电台便于携
行与机动使用，对功耗、质量与供电方式提出特殊要求；其三是\tex
tbf{可靠性}：须适应高低温、湿热、振动、盐雾等环境条件，关键
设备常按军标或工业级防护设计；其四是\textbf{灵活性}：支
持多种工作模式（话音、数据、加密）与多个频点，便于组网与应急转信
。

从网络视角看，电台是通信链路的两端实体。一条点对点链路的有效距离
取决于双方发射功率、天线增益、接收灵敏度与传播损耗；在组网应用中
，电台还需具备选呼、信道扫描、中继与网管能力。随着软件定义无线电
（SDR）的成熟，现代电台已能在同一硬件平台上通过加载不同波形软
件兼容多种体制，显著提升了互换性与生命周期成本优势。

电台的选型与配置须综合覆盖距离、业务密度与使用环境：远距离或存在
遮挡时优先考虑短波或中继链路；本地高密度调度宜选用超短波数字集群
；车载与便携设备强调电池续航与天线效率；固定基地则重视冗余电源与
高增益天线。与此同时，电台必须满足电磁兼容（EMC）与频谱管理规
定，其发射模板、杂散辐射与邻道泄漏均应达标，以免干扰同址部署的其
他无线电设备，这也是型号核准与入网检测的重点。

从节点属性看，电台既是通信网的末端也是中继节点：在点到点链路中它
独立完成收发，在网状/星状网络中它还需承担路由、转发与网管代理功
能。供电形态亦随平台不同而变化——手持台以锂聚合物电池供电并强调
低待机电流，车载台取用车载 $12/24\ \mathrm{V}
$ 电源并内置升压变换，基地台则采用交流整流与蓄电池浮充构成不间
断电源（UPS）。不同供电方式直接影响发射占空比与热设计，是工程
选型时不可忽视的约束。

\section{电台组成与原理}
\label{sec:radio_station_components}

电台的基本组成可划分为射频信道、基带与接口、控制以及电源天线四大
单元：

\begin{itemize}
    \item \textbf{发射部分}：调制器、上变频、
驱动与功率放大器、天线开关，完成基带信息到射频辐射的转换。
    \item \textbf{接收部分}：低噪声放大、下变
频、中频/基带放大、解调器，完成射频信号到信息的恢复。
    \item \textbf{控制部分}：频率合成器、微处
理器/FPGA、用户界面，负责频点管理、模式切换与协议处理。
    \item \textbf{电源部分}：交直流变换、电池
管理与稳压，为各单元提供干净稳定的工作电压。
    \item \textbf{天线系统}：天线、馈线、匹配
网络与天调，实现射频能量的高效辐射与接收。
\end{itemize}

电台的工作原理可简述为：发射时将语音或数据调制到高频载波并放大后
由天线辐射；接收时通过调谐选择特定频率信号，经放大与解调还原出原
始信息。在收发共用天线的设备中，天线调谐器（天调）用于使天线在工
作频带内呈现接近特性阻抗，降低电压驻波比（VSWR），从而减少馈
线损耗并保护功放。天调本质是一个可程控的阻抗匹配网络，常以 LC
 拓扑实现。

其目标是在工作频带内将天线端的复阻抗变换到馈线特性阻抗（常用 $
50\ \Omega$），使电压驻波比（VSWR）尽量接近 1。
当 $Z_L$ 为负载阻抗、$Z_0$ 为馈线阻抗时，反射系数与
驻波比的关系为
\begin{equation}
\Gamma = \frac{Z_L-Z_0}{Z_L+Z_0}, \qquad \mathrm{VSWR} = \frac{1+|\Gamma|}{1-|\Gamma|}.
\end{equation}
驻波过大不仅造成馈线功率反射、降低有效辐射，还会使功放输出端承受
高电压/高电流而触发保护甚至损坏，因此天调与馈线匹配是电台可靠工
作的基础。对于短波宽带天线，天调往往需在多个频点自动搜索最佳匹配
网络，匹配速度与带宽覆盖是关键技术指标。

电台的收发方式分为双工与半双工。\textbf{双工}指收发可同
时进行：频分双工（FDD）上下行占用不同频点，需双工器隔离；时分
双工（TDD）上下行占用同一频点但分时隙交换。\textbf{半
双工}指同一时刻只能收或发，由按键（PTT）控制收发切换，结构较
简单、功耗较低。接收灵敏度与噪声带宽的关系可表示为
\begin{equation}
P_{\min} = k T_0 B F,
\end{equation}
其中 $k$ 为玻尔兹曼常数，$T_0=290\ \mathrm
{K}$ 为参考温度，$B$ 为噪声带宽，$F$ 为噪声系数（倍
数），该式说明降低带宽与噪声系数可提升微弱信号接收能力。

在半双工设备中，收发共用一副天线需借助天线开关或环行器实现隔离，
切换瞬间须抑制发射泄漏以免烧毁接收前端；双工设备则依靠双工器在收
发频段间提供足够隔离（通常数十 dB 量级）。接收通道普遍设有自
动增益控制（AGC）以稳定不同场强下的输出幅度，并借助静噪（Sq
uelch）在无信号时关闭音频输出。对于宽带软件电台，本振泄漏、
镜像频率与中频泄漏是影响相邻接收通道性能的关键因素，需通过合理的
频率规划、镜像抑制混频与滤波加以规避。

基带与接口单元负责将话音、数据与信令转换为适合调制的形式。模拟话
音经压缩编码降低速率，数字业务则经过信道编码（如纠错、交织）与成
帧后送调制器；控制部分以微处理器或 FPGA 实现频点管理、扫描
、加密与协议处理，并通过显示屏/按键或远程网管完成人机交互。良好
的总线与时钟设计能保证多制式波形在同一硬件上稳定切换，是软件定义
电台（SDR）工程实现的核心难点之一。

\section{电台类型与分类}
\label{sec:radio_station_types}

电台可按功率、频段、用途与调制方式等多维度分类，不同维度交叉组合
形成丰富的产品谱系。

\begin{itemize}
    \item \textbf{按功率}：手持电台（通常 $
<5\ \mathrm{W}$）、车载电台（$25\sim 10
0\ \mathrm{W}$）、基地电台（$>100\ \mat
hrm{W}$）。功率等级直接决定通信距离与供电方式。
    \item \textbf{按频段}：短波（HF，$3\
sim 30\ \mathrm{MHz}$）靠天波/地波远程通信
；超短波（VHF/UHF，$30\sim 300/300\sim
 1000\ \mathrm{MHz}$）以视距为主，适合本地组
网；微波（GHz 级）用于大容量中继与干线。
    \item \textbf{按用途}：军用电台、民用电台
、业余电台，分别侧重保密/加固、成本/易用与多功能/可玩性。
    \item \textbf{按调制方式}：AM 电台、F
M 电台、SSB 电台、数字电台（如各类语音与数据波形）。
\end{itemize}

手持、车载与基地台除功率差异外，在供电、天线与人机接口上也有区别
：手持台以电池供电、配软天线，强调轻量与待机；车载台由车辆电源供
电、外接鞭状天线，兼顾功率与机动；基地台常接入市电并配定向/全向
高增益天线与冗余电源，追求稳定连续工作。短波电台多需天调与较长天
线以匹配低频频段，超短波电台则常用简易鞭状天线即可获得良好驻波。

从使用形态看，手持台强调轻量、长待机与一键操作，天线多为柔性鞭状
且增益有限，覆盖以视距近距为主；车载台借助车身地网与外置天线显著
提升效率，功率与通信距离居中；基地台接入市电并配以高塔或定向天线
、双路电源，可进行长时间连续发射并承担中继与网管角色。短波电台因
频率低、对应天线电尺寸大，通常配备天调与可切换的鞭状/斜拉天线，
以在不同频段均获得可用驻波；超短波与微波电台天线则更加紧凑，多直
接耦合馈线。

在频率规划层面，电台使用须遵循频道划分与指配：每个频道由标称频率
与占用带宽定义，相邻频道按既定间隔排列；为避免互调与邻道干扰，组
网时应避开三阶互调产物频点，并预留保护带。对于多台同址设备，还需
考虑收发隔离、杂散发射限值与合路/分路损耗，必要时采用腔体双工器
与滤波器。通信距离可用自由空间路径损耗近似估算：
\begin{equation}
L_{\mathrm{fs}} = 20\log_{10}(d) + 20\log_{10}(f) + 32.44,
\end{equation}
其中 $d$ 单位 km、$f$ 单位 MHz；该式说明频率越高
、距离越远，路径损耗增长越快，需相应提高功率或增益来补偿。

\begin{table}[htbp]
\centering
\caption{按频段划分的电台特性对照}
\label{tab:station_band}
\begin{tabular}{|p{2.6cm}|p{3.4cm}|p{4.2cm}|p{3.4cm}|}
\hline
\textbf{频段} & \textbf{典型传播} & \textbf{主要用途} & \textbf{通信距离} \\
\hline
短波 HF & 电离层反射 & 远程/跨国通信 & 数百至数千 km \\
\hline
超短波 VHF & 视距/绕射 & 本地调度、海事 & 数十 km 内 \\
\hline
超短波 UHF & 视距/穿透 & 城市、楼宇内 & 数至数十 km \\
\hline
微波 & 视距直线 & 中继、干线链路 & 数十 km（需直视） \\
\hline
\end{tabular}
\end{table}

\section{电台技术指标}
\label{sec:radio_station_specs}

电台性能由发射与接收两端指标共同表征，主要技术指标包括：

\begin{itemize}
    \item \textbf{频率范围}：设备可工作的频率
覆盖区间，常分段给出。
    \item \textbf{输出功率}：发射机额定射频功
率，通常以峰值包络功率（PEP）或平均功率标示。
    \item \textbf{接收灵敏度}：达到额定输出信
噪比所需的最小输入电平，常用 $\mu\mathrm{V}$ 或
 $\mathrm{dBm}$ 表示。
    \item \textbf{选择性}：区分相邻信道信号的
能力，取决于中频滤波器带宽与形状因数。
    \item \textbf{频率稳定度}：频率准确度与短
期/长期漂移，由基准源决定。
    \item \textbf{调制方式}：设备支持的调制/
波形类型与速率。
\end{itemize}

信道规划中，\textbf{信道间隔} $\Delta f$ 指
相邻频道标称频率之差，常见取值为 $12.5\ \mathrm{
kHz}$ 或 $25\ \mathrm{kHz}$；已调信号占
用带宽须小于等于间隔以免邻道干扰。频率稳定度常以相对值表示：
\begin{equation}
S = \frac{\Delta f}{f_0},
\end{equation}
例如晶体基准的稳定度可达 $10^{-6}$ 量级，温补/恒温晶
振可进一步改善至 $10^{-7}\sim 10^{-9}$。接
收选择性常借助中频滤波器的矩形系数 $K_{r}$ 评价：
\begin{equation}
K_r = \frac{B_{60\ \mathrm{dB}}}{B_{6\ \mathrm{dB}}},
\end{equation}
$K_r$ 越接近 1，滤波器边沿越陡、邻道抑制越好。工程上还需
关注互调抑制、阻塞与杂散响应等指标，它们共同决定复杂电磁环境下的
可用性。

除接收端指标外，发射端同样有严格规范：输出功率的准确度与稳定度、
调制频偏或调幅指数上限、发射机杂散与谐波抑制、邻道泄漏比（ACL
R）以及占用带宽均须符合无线电管理要求。在数据业务中，还需关注误
码率（BER）与吞吐随信干噪比（SINR）的变化关系。实际选型常
以"灵敏度—功率—选择性"三角权衡：提高功率可延长距离但增加干扰
与耗电，改善选择性可抗邻道却压缩带宽，须结合网络规划综合确定。对
于组网设备，还应考察越区切换时延、呼叫建立时间与多点接入容量，这
些指标共同决定实战中的通信可用性。

\section{电台应用与维护}
\label{sec:radio_station_maintenance}

电台凭借即时、不依赖公共网络的特点，在诸多领域不可替代：

\begin{itemize}
    \item \textbf{军事通信}：战术通信、指挥通
信，强调抗扰、保密与组网机动。
    \item \textbf{公共安全}：警察、消防、急救
通信，要求一键可达与越区切换。
    \item \textbf{交通运输}：航空、航海、铁路
通信，依国际/行业规则分配频点。
    \item \textbf{业余无线电}：爱好者通信、应
急通信，兼具技术实验属性。
\end{itemize}

为保障长期可靠工作，应建立周期性维护制度。常见维护项目包括：定期
通电老化与自检测试、用频率计/频谱仪校准载波频率、检查天线馈线与
接插件有无氧化或进水、测量驻波比确认天馈匹配、清洁散热风道并复核
电源电压纹波。一般手持/车载台建议每 $3\sim 6$ 个月做
一次功能检查，基地台则建议每月巡检并每年进行计量级指标复测；在潮
湿、盐雾或高尘环境中应适当缩短周期。

\begin{table}[htbp]
\centering
\caption{电台常见故障与维护周期参考}
\label{tab:station_maint}
\begin{tabular}{|p{3.4cm}|p{5.2cm}|p{3.0cm}|}
\hline
\textbf{常见故障现象} & \textbf{可能原因} & \textbf{维护周期} \\
\hline
发射功率下降 & 功放老化、馈线损耗增大 & 每 $3\sim 6$ 月 \\
\hline
接收灵敏度变差 & 天线失配、滤波器脏污 & 每 $3\sim 6$ 月 \\
\hline
频率漂移 & 晶振温漂、基准失效 & 每年计量 \\
\hline
通话断续 & 电池亏电、PTT 接触不良 & 每月巡检 \\
\hline
驻波比告警 & 天线损坏、接头进水 & 随告警即查 \\
\hline
\end{tabular}
\end{table}

维护时务必先断开天线再测量端口，避免大功率反向损坏仪表；更换功放
或天线后应重新测试驻波与发射模板。对于关键任务电台，建议配置冷备
份与备件，并定期演练应急开通流程，确保灾害或网络中断时通信不中断
。

现代电台普遍支持远程状态回传与健康监测，可将功放温度、驻波、电压
与发射功率等参数上报网管，实现预测性维护而非事后抢修。定期维护应
建立台账：记录每次校准的频率偏差、功率衰减趋势与电池循环次数，借
此判断器件老化拐点。一般建议手持/车载台每 $3\sim 6$ 
个月做一次功能检查，基地台每月巡检、每年计量复测；在潮湿、盐雾或
高粉尘环境中应缩短周期，并对接插件涂抹防氧化剂、对馈线接头做防水
处理。通过规范化维护，可显著延长电台寿命并降低任务中断概率。

在应急与野外部署中，电台常构成无中心（Ad-hoc）或稀疏骨干的
专用网络：通过中继台延伸视距、利用短波实现超视距备份，并以卫星或
北斗短报文回传关键报文。为降低全生命周期运维成本，应建立备件库与
故障树，对功放、电源与天线等易损件预置替换单元，并定期对操作手开
展开通与排障培训。多数行业（如公安、海事、航空、铁路）对电台频率
、呼号与制式有专门规定，部署前应完成频率申请与电磁环境测试，确保
合法合规、互不干扰。故障定位通常遵循"先电源、后天馈、再射频、后
基带"的顺序：先测供电与驻波，再逐级测量发射功率与接收灵敏度，最
后借助自环与仪表定位调制/解调与协议层问题，可显著提高修复效率。
